\section{A strict comparison of successive least dimensions}

The following auxiliary gives the dimension argument in
\cite[Proposition~prop:albootstrap-ld, lines 1112--1128]{Arad2016RigorousRG}
for nested targets and improving families. Its hypotheses describe the
families explicitly; their construction for a physical Hamiltonian is
separate. The sample size corrects the missing strict decrease in the
displayed source argument, as explained
in~\cite{gap:alvv16_low_density_sampling}.

\begin{theorem}[A bound from nested improving families]
    \label{thm:viable_dyadic_minimum_bound}
    \lean{Submodule.exists_isViable_finrank_le_of_dyadic_families}
    \leanok
    Let $T_j\subseteq T_{j+1}\subseteq M\otimes F$, where $M,F$ are
    finite-dimensional complex inner-product spaces. Let $r\ge1$ and
    $D_j$ be natural numbers such that $\dim T_j\le r$,
    $\log D_j\ge100$ and
    $\log D_{j+1}\le(13/6)\log D_j$.
    Suppose $\dim M\le rD_j^3$ for some finite index $j$.
    For each $j$, suppose there are $D_j^2$ regional operators $A_i$,
    targets $T_-,T_+$ and $\Delta_j\ge0$ such that $T_{j+1}$ is
    $1/200$-close to $T_+$, $T_-$ is $1/200$-close to $T_j$, and
    $D_j^{12}\Delta_j\le10^{-5}$. Suppose also that for every
    $\delta<1$ and every $\delta$-viable space $S$ for $T_+$,
    the space $\sum_i A_iS$ is
    $\Delta_j/(1-\delta)^2$-viable for $T_-$.
    Then there is a $3/200$-viable space $W$ for $T_0$ with
    $\dim W\le rD_0^3$.
\end{theorem}
\begin{proof}
    \leanok
    \uses{thm:viable_minimum_sequence,lem:nested_actual_sample_from_counts}
    Choose attained least dimensions $p_j$ and corresponding viable
    spaces $W_j$. The finite count bound gives a least $j$ with
    $p_j\le rD_j^3$. If $j=0$, use $W_0$.
    Otherwise write $j=k+1$. Then
    $rD_k^3<p_k\le p_{k+1}\le rD_{k+1}^3$.
    Draw the calibrated sample in $W_{k+1}$ with dimension
    $s=\lfloor p_k/(2D_k^2)\rfloor$.
    Its improving image is $3/200$-viable for $T_k$ and has dimension
    at most $p_k/2<p_k$, contradicting the minimum property of $p_k$.
    The sample dimension uses $p_k$, whereas the space in which it is
    sampled has dimension $p_{k+1}$.
\end{proof}

\begin{theorem}[The actual dyadic ceiling counts]
    \label{thm:viable_dyadic_filter_minimum_bound}
    \lean{Submodule.exists_isViable_finrank_le_of_dyadic_filter_families}
    \leanok
    Fix $C\ge100$, $v>0$ and an integer $m\ge2$ with $\log m\ge100$.
    Define $g_j,u_j,\ell_j,D_j$ as in
    Theorem~\ref{thm:dyadic_filter_count_finite_index}.
    Suppose $r\ge1$, the targets $T_j$ increase, $\dim T_j\le r$,
    and the nested improving families satisfy the closeness,
    suppression and improvement conditions of
    Theorem~\ref{thm:viable_dyadic_minimum_bound}.
    Then a $3/200$-viable space $W$ for $T_0$ satisfies
    $\dim W\le rD_0^3$. No count comparison or finite count bound
    is an additional hypothesis.
\end{theorem}
\begin{proof}
    \leanok
    \uses{thm:dyadic_filter_count_finite_index,
        thm:log_polynomial_filter_count_half_step,
        thm:viable_dyadic_minimum_bound}
    The identity $v/(6\cdot2^{j+1})=(v/(6\cdot2^j))/2$ gives
    $g_j\le2g_{j+1}$. Apply the dyadic count comparison.
    A finite index with $D_j\ge\dim M$ gives
    $\dim M\le D_j\le rD_j^3$, so the first theorem applies.
\end{proof}
